\documentclass[../../../include/open-logic-section]{subfiles}

\begin{document}

\olfileid[hi]{sth}{z}{union}
\olsection{सम्मिलन}

\olref[sth][z][sep]{prop:intersectionsexist} ने हमें सर्वनिष्ठ दिए। किंतु
यदि हम मनमाने सम्मिलनों का अस्तित्व चाहते हैं, तो एक और अभिगृहीत निर्धारित
करना होगा:

\begin{axiom}[सम्मिलन] किसी भी समुच्चय $A$ के लिए समुच्चय $\bigcup A =
\Setabs{x}{(\exists b \in A) x \in b}$ का अस्तित्व है।
\[
	\forall A \exists U \forall x(x \in U \liff (\exists b \in A)x \in b)
\]
\end{axiom}

इस अभिगृहीत का औचित्य भी संचयी-पुनरावर्ती अवधारणा से मिलता है। $A$ को
समुच्चय मानें, अतः $A$ किसी चरण $S$ पर बना है (\stageshier{} से)। $A$ का
प्रत्येक सदस्य $S$ से \emph{पहले} बना था (\stagesacc{} से); इसलिए समान
तर्क से $A$ के प्रत्येक सदस्य का प्रत्येक सदस्य $S$ से पहले बना था। अतः वे
\emph{सभी} समुच्चय $S$ से पहले उपलब्ध हैं ताकि $S$ पर उनसे एक समुच्चय बनाया
जा सके। और वह समुच्चय ठीक $\bigcup A$ है।

\end{document}
